\SourceLeaf{146}
though I hope you will allow me to keep it a few weeks longer. From a preliminary glance it seems to me to contain much that is admirable. I find in this book, too, much of what was distinctive in my method, especially in its earliest form. It seems to be my fate to compete with \textsc{Legendre} in almost all my theoretical work. So it was in higher arithmetic, in the investigations of transcendental functions connected with the rectification of the ellipse, in the first principles of geometry, and now once again here. Thus, e.\,g.\ the principle I have used since 1794, that to represent as well as possible several quantities that cannot all be represented exactly one must make the sum of the squares a minimum, is likewise adopted and put into practice quite ably in \textsc{Legendre}'s work.

\begin{center}
. . . . . . . . . . . . .
\end{center}

\begin{center}
\textsc{Gauss} to \textsc{Olbers}. Göttingen, 4 October 1809.
\end{center}

. . . . . . . . . Do you still remember, my dearest friend, that during my first visit to Bremen in 1803 I spoke to you about the principle I used to represent observations as accurately as possible, namely that where the observations are of equal weight the sum of the squares of the differences must be a minimum? I still remember every circumstance of our discussion of it in 1804 in Rehburg. It matters to me to know this. I shall tell you the reason for the question another time.

\begin{center}
. . . . . . . . . . . . . . .
\end{center}
